\section*{Bab \ref{ch:b3:lebesgue} --- Integral
Lebesgue}

\begin{solution}{exo:b3:lebesgue:1}
(a) Jika $f$ tak turun, maka $\{f > t\}$ berupa $\varnothing$, $\R$,
atau sinar $\intoo a{+\infty}$ / $\intco a{+\infty}$: yang bersifat Borel dalam
setiap kasusnya; dan yang tak naik serupa saja. Untuk turunan: $f'(x) =
\lim_n n\bigl(f(x + \frac1n) - f(x)\bigr)$ merupakan limit titik demi titik
fungsi \omterm{def:b3:topology:continuity}{kontinu} (jadi \omterm{def:b3:lebesgue:measurable}{terukur}):
\cref{prop:b3:lebesgue:stability}(c).

(b) $\{f > t\} = \bigcup_{q \in \Q,\, q > t}\{f > q\}$: sehingga jika
aras rasionalnya \omterm{def:b3:lebesgue:measurable}{terukur}, maka semua arasnya \omterm{def:b3:lebesgue:measurable}{terukur}, dan sinarnya
membangkitkan $\mathcal B(\R)$.
\end{solution}

\begin{solution}{exo:b3:lebesgue:2}
Pertama: $(1 + x/n)^n \nearrow \eu^x$ untuk $x \geq 0$, sehingga
integrannya menuju $\eu^{-x}\cos x$ secara titik demi titik; lalu untuk $n \geq 2$,
$(1 + x/n)^n \geq (1 + x/2)^2$, yang memberikan pendominasi \omterm{def:b3:lebesgue:l1}{terintegralkan}
$(1 + x/2)^{-2}$. Lewat kekonvergenan terdominasi:
\[
\lim_n\int_0^\infty\frac{\cos x}{(1 + x/n)^n}\dd x
= \int_0^\infty \eu^{-x}\cos x\,\dd x
= \operatorname{Re}\int_0^\infty\eu^{-(1 - \iu)x}\dd x
= \operatorname{Re}\frac{1}{1 - \iu} = \frac12 .
\]
Kedua: substitusikan $u = x^n$ (yakni bijeksi $\mathcal C^1$ dari
$\intoo01$):
\[
\int_0^1\frac{nx^{n-1}}{1 + x}\dd x = \int_0^1\frac{\dd u}{1 +
u^{1/n}} \longrightarrow \int_0^1\frac{\dd u}{2} = \frac12,
\]
lewat kekonvergenan terdominasi: sebab untuk $u \in \intoo01$, $u^{1/n} \to 1$, dan
integrannya terbatas oleh $1$ pada ruang berukuran berhingga.
\end{solution}

\begin{solution}{exo:b3:lebesgue:3}
(a) $f_n = n\,\mathbf 1_{\intoo0{1/n}}$: maka $\liminf f_n = 0$
secara titik demi titik, padahal $\int f_n = 1$: sehingga $0 < 1$.

(b) Secara tinggi: $n\mathbf 1_{\intoo0{1/n}}$; secara lebar: $\frac1n\mathbf
1_{\intoo0n}$; lewat translasi: $\mathbf 1_{\intoo n{n+1}}$. Semuanya
menuju $0$ secara titik demi titik dengan $\int = 1$. Dalam setiap kasusnya, hipotesis
\emph{dominasinya} gagal: sebab $\sup_nf_n$ bernilai
$\approx 1/x$ di dekat $0$, $\approx$ sebuah profil konstan yang tak
\omterm{def:b3:lebesgue:l1}{terintegralkan}, dan mirip $\mathbf 1_{\intoo1\infty}$ --- tak pernah \omterm{def:b3:lebesgue:l1}{terintegralkan}.

(c) Pernyataan ``$\int\limsup f_n \geq \limsup\int f_n$'' gagal bagi
gundukan yang bertranslasi: sebab ruas kirinya $0$, ruas kanannya $1$.
Sedangkan ``$\limsup\int \leq \int\limsup$'' adalah pernyataan yang sama. Dan
Fatou untuk $\limsup$ dengan $\leq$ yang terbalik
(yakni ``Fatou terbalik'') menuntut sebuah pendominasi --- dan gundukan yang sama
menjadi contoh penyangkalnya.
\end{solution}

\begin{solution}{exo:b3:lebesgue:4}
(a) Untuk $x > 0$: $\frac1{\eu^x - 1} = \frac{\eu^{-x}}{1 -
\eu^{-x}} = \sum_{n\geq1}\eu^{-nx}$, sehingga $\frac{x}{\eu^x - 1} =
\sum_{n\geq1}x\eu^{-nx}$, yakni deret berisi \omterm{def:b3:lebesgue:measurable}{fungsi terukur} tak
negatif: sehingga \cref{cor:b3:lebesgue:additivity} mengizinkan pengintegralan
suku demi suku:
\[
\int_0^\infty\frac{x\,\dd x}{\eu^x - 1}
= \sum_{n\geq1}\int_0^\infty x\eu^{-nx}\dd x
= \sum_{n\geq1}\frac1{n^2} = \frac{\pi^2}6
\]
(sebab $\int_0^\infty x\eu^{-nx}\dd x = n^{-2}$ lewat pengintegralan parsial; sedangkan Basel dari
jilid Tahun ke-2, atau \cref{exo:b3:hilbert:5} yang akan datang).

(b) Pada $\intoo01$, $-x\ln x \geq 0$, sehingga $x^{-x} =
\eu^{-x\ln x} = \sum_k\frac{(-x\ln x)^k}{k!}$ merupakan deret
bersuku tak negatif: jadi tukarkanlah lagi. Dengan substitusi $x =
\eu^{-u/(k+1)}$:
\[
\int_0^1(-x\ln x)^k\dd x
= \int_0^\infty\Bigl(\frac{u}{k+1}\Bigr)^{k}
\eu^{-\frac{ku}{k+1}}\;\frac{\eu^{-\frac u{k+1}}}{k+1}\,\dd u
= \frac{1}{(k+1)^{k+1}}\int_0^\infty u^k\eu^{-u}\dd u
= \frac{k!}{(k+1)^{k+1}} .
\]
Karena itu $\int_0^1x^{-x}\dd x = \sum_{k\geq0}\frac{1}{(k+1)^{k+1}}
= \sum_{n\geq1}n^{-n}$: yakni mimpi mahasiswa tahun kedua, secara ketat.
\end{solution}

\begin{solution}{exo:b3:lebesgue:5}
(a) Markov: $a\,\mathbf 1_{\{f \geq a\}} \leq f$, lalu integralkan:
$\mu(\{f \geq a\}) \leq \frac1a\int f$. Jika $\int f = 0$: maka
$\mu(\{f \geq \frac1n\}) = 0$ untuk setiap $n$, dan $\{f > 0\} =
\bigcup_n\{f \geq \frac1n\}$ bersifat nol.

(b) $\mu(\{\abs f = \infty\}) \leq \mu(\{\abs f \geq n\}) \leq
\frac1n\int\abs f \to 0$.

(c) Ambillah $A = \{f > 0\}$: maka $\int f^+\dd\mu = \int_Af\,\dd\mu =
0$, sehingga $f^+ = 0$ \omterm{def:b3:lebesgue:l1}{hampir di mana-mana} menurut (a); dan demikian pula $f^- = 0$.
\end{solution}

\begin{solution}{exo:b3:lebesgue:6}
(a) Fungsi $\mathbf 1_\Q$: tak \omterm{def:b3:topology:continuity}{kontinu} di mana pun, sehingga tak
\omterm{def:b3:lebesgue:l1}{terintegralkan} Riemann (\cref{thm:b3:lebesgue:riemann}); sedangkan
integral Lebesguenya $\lambda(\Q) = 0$. Fungsi Thomae: \omterm{def:b3:topology:continuity}{kontinu} di
setiap bilangan irasional (sebab diberikan $\varepsilon$, hanya berhingga banyak
bilangan rasional di $\intcc01$ yang berpenyebut $\leq
1/\varepsilon$; lalu hindarilah semuanya lewat \omterm{def:b3:topology:topology}{persekitaran} kecil),
dan tak \omterm{def:b3:topology:continuity}{kontinu} di bilangan rasional (menurut \omterm{ex:b3:lebesgue:gamma}{kepadatan} bilangan irasional): jadi \omterm{def:b3:topology:continuity}{kontinu}
\omterm{def:b3:lebesgue:l1}{hampir di mana-mana}, \omterm{def:b3:lebesgue:l1}{terintegralkan} Riemann, dengan integral $0$ (sebab ia lenyap \omterm{def:b3:lebesgue:l1}{hampir di mana-mana}).
Sedangkan $\mathbf 1_K$ dengan $K$ himpunan Cantor gemuk: himpunan ketakkontinuannya
adalah $\partial K = K$ (yang tertutup dengan \omterm{def:b3:topology:interior}{interior} kosong), berukuran
$\frac12 > 0$: jadi \emph{tak} \omterm{def:b3:lebesgue:l1}{terintegralkan} Riemann --- padahal
\omterm{def:b3:lebesgue:l1}{terintegralkan} Lebesgue dengan integral $\lambda(K) = \frac12$.

(b) $\int_{k\pi}^{(k+1)\pi}\frac{\abs{\sin x}}x\dd x \geq
\frac1{(k+1)\pi}\int_{k\pi}^{(k+1)\pi}\abs{\sin x}\dd x =
\frac{2}{(k+1)\pi}$: sehingga deretnya divergen. Untuk kekonvergenan
integral tak wajarnya: bila $A > \pi$,
\[
\int_\pi^A\frac{\sin x}x\dd x = \Bigl[\frac{-\cos
x}x\Bigr]_\pi^A - \int_\pi^A\frac{\cos x}{x^2}\dd x,
\]
dan kedua sukunya konvergen saat $A \to \infty$ (sebab $\frac1{x^2}$
\omterm{def:b3:lebesgue:l1}{terintegralkan}): yakni kekonvergenan bersyarat tanpa keterintegralan
mutlak.
\end{solution}

\begin{solution}{exo:b3:lebesgue:7}
Dominasinya: $\abs{\partial_t(\eu^{-x^2}\cos(tx))} =
\abs{x\eu^{-x^2}\sin(tx)} \leq x\eu^{-x^2}$, yang \omterm{def:b3:lebesgue:l1}{terintegralkan} dan
tak bergantung pada $t$: sehingga \cref{thm:b3:lebesgue:paramdiff} berlaku
secara global,
\[
F'(t) = -\int_0^\infty x\eu^{-x^2}\sin(tx)\,\dd x
= \Bigl[\tfrac12\eu^{-x^2}\sin(tx)\Bigr]_0^\infty
- \frac t2\int_0^\infty\eu^{-x^2}\cos(tx)\dd x
= -\frac t2F(t)
\]
(lewat pengintegralan parsial dengan $\dd v = x\eu^{-x^2}\dd x$). Lalu
persamaan diferensial linearnya memberikan $F(t) = F(0)\eu^{-t^2/4}$; dan dengan $F(0) =
\frac{\sqrt\pi}2$ (\cref{pb:b3:lebesgue:1}), fungsi Gauss
mereproduksi dirinya di bawah transformasi kosinus ini.
\end{solution}

\begin{solution}{exo:b3:lebesgue:8}
Integralnya konvergen: sebab di dekat $0$ integrannya menuju $b - a$
(yang terbatas), dan ia meluruh seperti $\eu^{-ax}$ di tak hingga. Tetapkan $a$
lalu pandanglah $I(b) = \int_0^\infty\frac{\eu^{-ax} -
\eu^{-bx}}x\dd x$ sebagai fungsi $b \in \intco a{+\infty}$.
Untuk setiap $b \geq a$: $\abs{\partial_b(\text{integran})} =
\eu^{-bx} \leq \eu^{-ax}$, dan $\int_0^\infty\eu^{-ax}\dd x =
\frac1a < \infty$: yakni pendominasi yang \omterm{def:b3:lebesgue:l1}{terintegralkan}. Sehingga
\cref{thm:b3:lebesgue:paramdiff} memberikan $I'(b) =
\int_0^\infty\eu^{-bx}\dd x = \frac1b$, dengan $I(a) = 0$:
\[
I(b) = \int_a^b\frac{\dd t}t = \ln\frac ba .
\]
(Setara dengan itu, lewat jalur petunjuknya: integrannya adalah
$\int_a^b\eu^{-xt}\dd t \geq 0$ dan penukarannya merupakan
analogi \omterm{def:b3:topology:continuity}{kontinu} \cref{cor:b3:lebesgue:additivity},
yakni Tonelli --- yang dibuktikan di \cref{ch:b3:product}; sedangkan
jalur parameternya tetap berada di dalam bab ini.)
\end{solution}

\begin{solution}{exo:b3:lebesgue:9}
Berlaku $\nu(\varnothing) = 0$; lalu untuk $(A_n)$ yang saling lepas, $f\mathbf
1_{\bigsqcup A_n} = \sum_nf\mathbf 1_{A_n}$ (secara titik demi titik, dengan semua
sukunya $\geq 0$), dan \cref{cor:b3:lebesgue:additivity} memberikan
keaditifan-$\sigma$-nya. Untuk rumus $\int g\,\dd\nu = \int
gf\,\dd\mu$: bila $g = \mathbf 1_A$ ia adalah definisi
$\nu$; bila $g$ sederhana, lewat kelinearannya; sedangkan bila $g \geq 0$ \omterm{def:b3:lebesgue:measurable}{terukur},
ambillah $s_n \nearrow g$ yang sederhana
(\cref{thm:b3:lebesgue:approximation}): maka $s_nf \nearrow gf$, dan
kekonvergenan monoton pada kedua ruasnya mengalihkannya ke limitnya. (Eskalator ``indikator
$\to$ sederhana $\to$ kekonvergenan monoton'' inilah mesin baku
teorinya.)
\end{solution}

\begin{solution}{exo:b3:lebesgue:10}
(a) Dari $\abs{\sin(tx)} \leq \abs tx$ diperoleh
$\abs{\frac{\sin(tx)}{x(1+x^2)}} \leq \frac{\abs
t}{1+x^2}$: sehingga integralnya konvergen, dan pada $\abs t \leq T$
pendominasi $\frac{T}{1+x^2}$ melahirkan \omterm{def:b3:topology:continuity}{kekontinuannya}
(\cref{thm:b3:lebesgue:paramcont}). Untuk pendiferensialannya:
$\abs{\partial_t} = \abs{\frac{\cos(tx)}{1+x^2}} \leq
\frac1{1+x^2}$, yang \omterm{def:b3:lebesgue:l1}{terintegralkan}: sehingga $f'(t) =
\int_0^\infty\frac{\cos(tx)}{1+x^2}\dd x$ untuk setiap $t$. Sedangkan pendiferensialan
kedua akan menuntut pengintegralan $\frac{x\sin(tx)}{1 +
x^2}$, yang nilai mutlaknya berperilaku seperti $\frac{\abs{\sin(tx)}}
x$ di tak hingga: jadi \emph{tak} \omterm{def:b3:lebesgue:l1}{terintegralkan} --- sehingga tak ada pendominasi
dan \cref{thm:b3:lebesgue:paramdiff} tak dapat diterapkan lagi.

(b) Dengan menerima $f'(t) = \frac\pi2\eu^{-t}$ untuk $t > 0$: menurut
kegasalan $f$, $f'$ bersifat genap, sehingga $f'(t) =
\frac\pi2\eu^{-\abs t}$ --- yang \emph{tak dapat diturunkan
di $0$}: jadi $f$ bersifat $\mathcal C^1$ tetapi bukan $\mathcal C^2$,
yang membenarkan bahwa pendiferensialan kedua yang terhalang itu bukan
kecelakaan teknis. Untuk $t > 0$ integral tak wajar
$\int_0^\infty\frac{x\sin(tx)}{1+x^2}\dd x$ sama dengan $-f''(t) =
\frac\pi2\eu^{-t}$ (dengan menurunkan rumus yang diterima itu di tempat
yang sah, yakni pada $\intoo0\infty$) --- yaitu nilai tak wajar
yang bukan nilai Lebesgue.
\end{solution}

\begin{solution}{exo:b3:lebesgue:11}
(a) Misalkan $g_n = (f - f_n)^+$: maka $0 \leq g_n \leq f$
(menurut kepositifan $f_n$), $g_n \to 0$ \omterm{def:b3:lebesgue:l1}{hampir di mana-mana}, dan $f$ merupakan
pendominasi yang \omterm{def:b3:lebesgue:l1}{terintegralkan}: sehingga $\int g_n \to 0$ (lewat kekonvergenan terdominasi). Karena
$\abs{f_n - f} = 2(f - f_n)^+ - (f - f_n)$, maka
\[
\int\abs{f_n - f} = 2\int g_n - \Bigl(\int f - \int
f_n\Bigr) \longrightarrow 0 + 0 .
\]

(b) Gundukan yang meluncur $f_n = \mathbf 1_{\intcc n{n+1}}$ mempunyai
$f_n \to 0$ \omterm{def:b3:lebesgue:l1}{hampir di mana-mana} dan $\int f_n = 1 \not\to 0$: jadi tanpa
kekonvergenan integralnya, kekonvergenan $L^1$-nya gagal (dan begitu pula
hipotesisnya). Untuk contoh bertandanya: $f_n \to 0$ di
setiap $x \neq 0$, $\int f_n = 0$, tetapi $\int\abs{f_n} = 2$:
sehingga bagi barisan bertanda teoremanya sungguh berkenaan dengan
kendali bertipe $\abs{f_n}$, dan kepositifannya dipakai persis pada
$g_n \leq f$.

(c) Kerapatan memenuhi $\int p_n = 1 = \int p$: sehingga
hipotesis (a) berlaku dengan sendirinya, jadi $p_n \to p$ \omterm{def:b3:lebesgue:l1}{hampir di mana-mana} memaksa
$\norm{p_n - p}_{L^1} \to 0$ --- dan karena itu kekonvergenan
peluang $\int_Ap_n \to \int_Ap$ berlaku \emph{secara seragam}
atas semua $A$ yang \omterm{def:b3:lebesgue:measurable}{terukur} (sebab $\abs{\int_A(p_n - p)} \leq
\norm{p_n - p}_1$): jadi Scheffé mengubah kekonvergenan titik demi titik
\omterm{ex:b3:lebesgue:gamma}{kepadatan} menjadi kekonvergenan variasi total hukumnya.
\end{solution}

\begin{solution}{exo:b3:lebesgue:12}
(a) Pada $\intoo0n$, $\varphi_n(x) = n\log(1 - \frac xn)$
naik terhadap $n$ menuju $-x$ (sebab pemetaan $t \mapsto
\frac{\log(1 - xt)}{t}$ turun saat $t = \frac1n
\downarrow 0$; atau uraikanlah: $\varphi_{n}' \geq 0$ terhadap $n$ lewat
$\log(1-u) + \frac{u}{1-u} \geq 0$). Jadi $(1 -
\frac xn)^n\eu^{x/2}\mathbf 1_{x<n} \nearrow
\eu^{-x/2}$, dan kekonvergenan monotonnya memberikan
\[
\lim_n\int_0^n\Bigl(1 - \frac xn\Bigr)^n\eu^{x/2}\dd x =
\int_0^\infty\eu^{-x/2}\dd x = 2 .
\]

(b) Substitusikan $u = x^n$ (dengan $x = u^{1/n}$ dan $n x^{n-1}\dd x =
\dd u$):
\[
\int_0^1\frac{nx^{n-1}}{1 + x}\dd x =
\int_0^1\frac{\dd u}{1 + u^{1/n}} .
\]
Untuk $u \in \intoo01$: $u^{1/n} \to 1$, sehingga integrannya
menuju $\frac12$, dan terdominasi oleh $1$: jadi limitnya
$\frac12$ (lewat kekonvergenan terdominasi). (Massa $nx^{n-1}$ memusat di $x
= 1$, tempat $\frac1{1+x} = \frac12$: dan substitusinya membuat
pemusatan itu terlihat.)

(c) Secara titik demi titik, $(1 + x/n)^n \nearrow \eu^x$ dan $x^{1/n} \to
1$ (untuk $x > 0$): sehingga integrannya menuju $\eu^{-x}$. Pendominasi
untuk $n \geq 2$: pada $\intoc01$, $x^{-1/n} \leq x^{-1/2}$ dan
$(1 + x/n)^{-n} \leq 1$: sehingga batasnya $x^{-1/2}$, yang \omterm{def:b3:lebesgue:l1}{terintegralkan}; sedangkan pada
$\intoo1\infty$, $x^{-1/n} \leq 1$ dan $(1 + x/n)^n \geq 1 +
\binom n2\frac{x^2}{n^2} \geq 1 + \frac{x^2}4$: sehingga batasnya
$\frac{4}{4 + x^2}$, yang \omterm{def:b3:lebesgue:l1}{terintegralkan}. Lewat kekonvergenan terdominasi:
\[
\lim_n\int_0^\infty\frac{\dd x}{(1 + x/n)^nx^{1/n}} =
\int_0^\infty\eu^{-x}\dd x = 1 .
\]
\end{solution}

\begin{solution}{pb:b3:lebesgue:1}
\textbf{1.} Fungsi $A$ bersifat $\mathcal C^1$ menurut teorema fundamental
kalkulus dan aturan rantai: $A'(t) =
2\eu^{-t^2}\int_0^t\eu^{-x^2}\dd x$. Untuk $B$:
$\partial_t\bigl(\frac{\eu^{-t^2(1+x^2)}}{1+x^2}\bigr) =
-2t\,\eu^{-t^2(1+x^2)}$, yang \omterm{def:b3:topology:continuity}{kontinu} dan terbatas pada $[t_0, T]
\times \intcc01$ untuk sembarang $0 < t_0 < T$ (sebab dominasi terbatas pada
ruang berukuran berhingga sudah cukup): jadi $B$ bersifat $\mathcal C^1$ pada
$\intoo0{+\infty}$ dengan
\[
B'(t) = -2t\int_0^1 \eu^{-t^2(1+x^2)}\dd x
= -2\eu^{-t^2}\int_0^1 t\,\eu^{-(tx)^2}\dd x
= -2\eu^{-t^2}\int_0^t\eu^{-u^2}\dd u = -A'(t)
\]
(lewat substitusi $u = tx$).

\textbf{2.} Berlaku $A(0) = 0$ dan $B(0) = \int_0^1\frac{\dd x}{1+x^2}
= \frac\pi4$. Lalu $A + B$ berturunan nol pada
$\intoo0{+\infty}$ dan \omterm{def:b3:topology:continuity}{kontinu} di $0$ ($B$ lewat
dominasi $\frac1{1+x^2}$ dan
\cref{thm:b3:lebesgue:paramcont}): sehingga $A + B \equiv \frac\pi4$.
Saat $t\to+\infty$: $0 \leq B(t) \leq \eu^{-t^2} \to 0$, dan
$A(t) \to \bigl(\int_0^\infty\eu^{-x^2}\dd x\bigr)^2$ (lewat kekonvergenan monoton atau
sekadar kekonvergenan monoton integral dalamnya).

\textbf{3.} Karena itu $\bigl(\int_0^\infty\eu^{-x^2}\dd x\bigr)^2
= \frac\pi4$: sehingga $\int_0^\infty\eu^{-x^2}\dd x =
\frac{\sqrt\pi}2$, dan menurut kegenapannya $G = \sqrt\pi$. Selain itu
$\Gamma(\frac12) = \int_0^\infty x^{-1/2}\eu^{-x}\dd x
\overset{x = u^2}{=} 2\int_0^\infty\eu^{-u^2}\dd u =
\sqrt\pi$.

\textbf{4.} Untuk $t > 0$: $\abs{\eu^{-tx}\frac{\sin x}x} \leq
\eu^{-tx}$, yang \omterm{def:b3:lebesgue:l1}{terintegralkan}. Untuk $t = 0$, kekonvergenan tak wajarnya: pada
$[\pi, A]$,
\[
\int_\pi^A\frac{\sin x}x\dd x
= \Bigl[-\frac{\cos x}x\Bigr]_\pi^A -
\int_\pi^A\frac{\cos x}{x^2}\dd x,
\]
dengan kedua sukunya konvergen saat $A \to \infty$; sedangkan di dekat $0$ integrannya
diperluas secara \omterm{def:b3:topology:continuity}{kontinu} oleh $1$.

\textbf{5.} Pada $[t_0, +\infty)$ (dengan $t_0 > 0$):
$\abs{\partial_t\bigl(\eu^{-tx}\tfrac{\sin x}x\bigr)} =
\eu^{-tx}\abs{\sin x} \leq \eu^{-t_0x}$, yang \omterm{def:b3:lebesgue:l1}{terintegralkan}: sehingga
\cref{thm:b3:lebesgue:paramdiff} berlaku pada setiap selang semacam
itu, jadi pada seluruh $\intoo0{+\infty}$:
\[
F'(t) = -\int_0^\infty\eu^{-tx}\sin x\,\dd x
= -\operatorname{Im}\int_0^\infty\eu^{-(t - \iu)x}\dd x
= -\operatorname{Im}\frac{1}{t - \iu} = -\frac{1}{1 + t^2}.
\]

\textbf{6.} Berlaku $\abs{F(t)} \leq \int_0^\infty\eu^{-tx}\dd x =
\frac1t \to 0$. Lalu dengan mengintegralkan $F' = -\frac1{1+t^2}$: $F(t) = C -
\arctan t$, dan $t \to \infty$ memaksa $C = \frac\pi2$: sehingga $F(t) =
\frac\pi2 - \arctan t$ pada $\intoo0{+\infty}$.

\textbf{7.} Integralkan secara parsial pada $[A, R]$ dengan $\sin x =
(-\cos x)'$ lalu biarkan $R \to \infty$:
\[
\int_A^{\infty}\eu^{-tx}\frac{\sin x}x\dd x
= \frac{\cos A\;\eu^{-tA}}{A}
- \int_A^\infty \cos x\,\Bigl(\frac tx +
\frac1{x^2}\Bigr)\eu^{-tx}\dd x .
\]
Dengan membatasi $\abs{\cos} \leq 1$: suku pertamanya $\leq \frac1A$;
sedangkan integralnya paling banyak $\int_A^\infty t\eu^{-tx}\frac{\dd x}x
+ \int_A^\infty\frac{\dd x}{x^2} \leq \frac1A\int_0^\infty
t\eu^{-tx}\dd x + \frac1A = \frac2A$. Totalnya: $\leq \frac 3A$,
secara seragam untuk $t \in [0, 1]$ (termasuk kasus $t = 0$). Kini
\[
\abs{F(t) - D} \leq
\Bigl|\int_0^A(\eu^{-tx} - 1)\,\frac{\sin x}x\,\dd x\Bigr| +
\frac6A .
\]
Pada $[0, A]$: $\abs{\eu^{-tx} - 1} \leq tx \leq tA$ dan
$\abs{\frac{\sin x}x} \leq 1$, sehingga suku pertamanya paling banyak
$tA^2$. Pilihlah $A$ dengan $\frac6A < \varepsilon$, lalu $t <
\varepsilon/A^2$: sehingga $\abs{F(t) - D} < 2\varepsilon$.

\textbf{8.} Karena itu $D = \lim_{t\to0^+}F(t) =
\lim_{t\to0^+}\bigl(\frac\pi2 - \arctan t\bigr) =
\frac\pi2$.

\textbf{9.} Lewat pengintegralan parsial (dengan $u = \sin^2x$ dan $v' = x^{-2}$):
\[
\int_0^\infty\frac{\sin^2x}{x^2}\dd x
= \Bigl[-\frac{\sin^2x}{x}\Bigr]_0^\infty +
\int_0^\infty\frac{2\sin x\cos x}{x}\dd x
= \int_0^\infty\frac{\sin 2x}{x}\dd x = D = \frac\pi2
\]
(substitusikan $u = 2x$ pada langkah terakhirnya; sedangkan suku batasnya lenyap:
sebab $\sin^2 x/x \to 0$ di kedua ujungnya).

\textbf{10.} Lewat pengintegralan parsial (dengan $u = 1 - \cos x$ dan $v' = x^{-2}$):
$\int_0^\infty\frac{1 - \cos x}{x^2}\dd x =
\int_0^\infty\frac{\sin x}x\dd x = \frac\pi2$. Kekonsistenannya:
$1 - \cos x = 2\sin^2\frac x2$, dan substitusi $x = 2u$
mengubah $\int\frac{2\sin^2(x/2)}{x^2}\dd x$ menjadi
$\int\frac{\sin^2u}{u^2}\dd u$: sehingga kedua perhitungannya bersesuaian.

\textbf{11.} Berlaku $\int_0^\infty\frac{\sin(ax)}x\dd x = \frac\pi2$
untuk setiap $a > 0$ (substitusikan $u = ax$: sebab integralnya
invarian terhadap penskalaan); sedangkan $\int_\R\eu^{-ax^2}\dd x = \sqrt{\pi/a}$
(substitusikan $u = \sqrt a\,x$). Jadi penskalaan pada argumennya membiarkan
integral Dirichletnya tetap dan membagi integral Gaussnya dengan
$\sqrt a$.

\textbf{12.} Sebuah pendominasi yang berlaku untuk setiap $t \in [0,1]$ harus
mendominasi $\sup_{t\in[0,1]}\abs{\eu^{-tx}\frac{\sin
x}x} = \abs{\frac{\sin x}x}$, yang tak \omterm{def:b3:lebesgue:l1}{terintegralkan}
(\cref{exo:b3:lebesgue:6}): sehingga kekonvergenan terdominasinya tak dapat menyeberangi $t = 0$. Sedangkan
pengganti pada pertanyaan 7 --- yakni ekor yang seragam kecil terhadap
parameternya, dengan bagian \omterm{def:b3:topology:compact}{kompaknya} ditangani kekonvergenan terdominasi --- merupakan pola
``keterintegralan seragam'' yang baku, dan ia muncul kembali setiap kali
kekonvergenan bersyarat berjumpa penukaran limit.

\textbf{13.} Jika $g = \log f$ cembung maka $f = \exp\circ g$
cembung (sebab eksponensial cembung dan naik: $f(\lambda x + (1-\lambda)
y) \leq \eu^{\lambda g(x) + (1-\lambda)g(y)} \leq \lambda f(x)
+ (1-\lambda)f(y)$, dengan langkah terakhirnya lewat kecembungan eksponensial di antara
titik $g(x), g(y)$). Untuk hasil kali dan substitusi afinnya:
logaritmanya mengubahnya menjadi jumlah dan substitusi afin
fungsi cembung. Hölder (dengan $\frac1p + \frac1q = 1$): untuk
$\int\abs u^p = \int\abs v^q = 1$, Young memberikan $\abs{uv} \leq
\frac{\abs u^p}p + \frac{\abs v^q}q$, lalu integralkan: $\int\abs{uv}
\leq 1$; sedangkan kasus umumnya lewat kehomogenannya. Lalu untuk $\lambda
\in \intoo01$, terapkanlah dengan $p = \frac1\lambda$ pada
pemfaktoran
\[
t^{\lambda x + (1-\lambda)y - 1}\eu^{-t}
= \bigl(t^{x-1}\eu^{-t}\bigr)^{\lambda}
\bigl(t^{y-1}\eu^{-t}\bigr)^{1-\lambda} :
\qquad
\Gamma(\lambda x + (1{-}\lambda)y) \leq
\Gamma(x)^\lambda\,\Gamma(y)^{1-\lambda} .
\]

\textbf{14.} Untuk $g$ yang cembung, kemiringan sebuah talibusur naik
seiring ujungnya (yakni ketaksamaan tiga talibusur). Dengan membandingkan
talibusur sepanjang $[n-1, n]$, $[n, n+x]$, dan $[n, n+1]$:
\[
\log(n-1) = \frac{g(n) - g(n-1)}1 \leq
\frac{g(n+x) - g(n)}x \leq \frac{g(n+1) - g(n)}1 = \log n,
\]
dan mengalikannya dengan $x > 0$ memberikan pernyataannya.

\textbf{15.} Dengan $f(n) = (n-1)!$ (yakni rekursi dari $f(1) = 1$)
dan $f(n + x) = x(x+1)\cdots(x+n-1)\,f(x)$, pertanyaan 14 berbunyi
\[
(n-1)^x\,(n-1)! \;\leq\; x(x+1)\cdots(x+n-1)\,f(x)
\;\leq\; n^x\,(n-1)! .
\]
Batas atasnya ditulis ulang sebagai $f(x) \leq
\frac{n!\,n^x}{x(x+1)\cdots(x+n)}\cdot\frac{x+n}n$, sedangkan
batas bawahnya pada peringkat $n+1$ sebagai $f(x) \geq
\frac{n!\,n^x}{x(x+1)\cdots(x+n)}$. Faktor koreksinya
$\frac{x+n}n \to 1$: sehingga penjepitannya memaksa
\[
f(x) = \lim_n\frac{n!\,n^x}{x(x+1)\cdots(x+n)}
\qquad (x \in \intoc01),
\]
yakni ungkapan yang tak bergantung pada $f$: jadi ketunggalannya pada $\intoc01$,
sehingga di mana-mana lewat rekursinya. Dan karena $\Gamma$ memenuhi
ketiga hipotesisnya (pertanyaan 13), maka $f = \Gamma$ dan rumus
Gauss berlaku --- untuk setiap $x > 0$, sebab kedua ruasnya menuruti
rekursi yang sama.

\textbf{16.} Di dekat $0$, integrannya $\sim t^{x-1}$, yang
\omterm{def:b3:lebesgue:l1}{terintegralkan} jika dan hanya jika $x > 0$; sedangkan di dekat $1$, simetris dengan $y$.
Lakukan pengintegralan parsial pada $\intcc\varepsilon{1-\varepsilon}$,
lalu biarkan $\varepsilon \to 0$ (suku batasnya lenyap untuk $x, y >
0$):
\[
B(x{+}1, y) = \Bigl[-t^x\frac{(1-t)^y}y\Bigr]_0^1 +
\frac xy\int_0^1t^{x-1}(1-t)^y\,\dd t
= \frac xy\bigl(B(x, y) - B(x{+}1, y)\bigr),
\]
dengan memakai $(1-t)^y = (1-t)^{y-1}(1 - t)$; lalu selesaikan: $B(x+1, y) =
\frac{x}{x+y}B(x, y)$. Dan $B(1, y) =
\int_0^1(1-t)^{y-1}\dd t = \frac1y$.

\textbf{17.} Berlaku $f(1) = \frac1y\cdot\frac{\Gamma(1+y)}{\Gamma(y)}
= 1$; lalu $f(x+1) = \frac{x}{x+y}B(x,y)\cdot\frac{(x+y)\Gamma(x+y)}
{\Gamma(y)} = x\,f(x)$; dan $f$ bersifat log-cembung terhadap $x$ sebagai
hasil kali $B(\cdot, y)$ yang log-cembung (lewat Hölder pada
pemfaktoran $t^{(\lambda x_1 + (1-\lambda)x_2)-1}(1-t)^{y-1}
= (\cdots)^\lambda(\cdots)^{1-\lambda}$, seperti pada pertanyaan 13)
dan $\Gamma(\cdot + y)$ (lewat pergeseran afin). Menurut Bohr--Mollerup: $f =
\Gamma$, yakni $B(x, y) =
\frac{\Gamma(x)\Gamma(y)}{\Gamma(x+y)}$.

\textbf{18.} Dengan $t = \sin^2\theta$, $\dd t =
2\sin\theta\cos\theta\,\dd\theta$ dan $t^{-1/2}(1-t)^{-1/2} =
\frac1{\sin\theta\cos\theta}$:
\[
B\Bigl(\frac12, \frac12\Bigr) =
\int_0^{\pi/2}2\,\dd\theta = \pi
= \frac{\Gamma(\frac12)^2}{\Gamma(1)}
\quad\Longrightarrow\quad
\Gamma\Bigl(\frac12\Bigr) = \sqrt\pi .
\]
Bagian I mencapai konstanta yang sama dengan menurunkan sebuah
parameter dan mengadu dua fungsi menuju limitnya; sedangkan di sini
kecembungan belaka mengakukan soalnya sampai hanya satu nilai yang
tersisa. Yakni analisis lewat gerak lawan analisis lewat bentuk.

\textbf{19.} Berlaku $g(1) = \frac{2^0}{\sqrt\pi}\Gamma(\frac12)
\Gamma(1) = 1$. Rekursinya:
\[
g(x+1) = \frac{2^{x}}{\sqrt\pi}\,
\Gamma\Bigl(\frac{x+1}2\Bigr)\Gamma\Bigl(\frac x2 + 1\Bigr)
= \frac{2^{x}}{\sqrt\pi}\cdot\frac x2\,
\Gamma\Bigl(\frac x2\Bigr)\Gamma\Bigl(\frac{x+1}2\Bigr)
= x\,g(x) .
\]
Kelog-cembungannya: yakni hasil kali $\eu^{(x-1)\log2}$ (yang log-afin) dan
dua reparametrisasi afin dari $\Gamma$ yang log-cembung.
Lalu Bohr--Mollerup memberikan $g = \Gamma$; dan dengan menetapkan $x = 2z$:
$\Gamma(2z) = \frac{2^{2z-1}}{\sqrt\pi}\Gamma(z)\Gamma(z +
\frac12)$.

\textbf{20.} Dari $\Gamma(\frac12) = \sqrt\pi$ dan
rekursinya, $\Gamma(n + \frac12) = (n - \frac12)(n -
\frac32)\cdots\frac12\,\sqrt\pi =
\frac{(2n-1)(2n-3)\cdots1}{2^n}\sqrt\pi =
\frac{(2n)!}{4^nn!}\sqrt\pi$ (lengkapilah hasil kali gasalnya dengan
yang genap). Keasimtotikannya: pertanyaan 14 dengan $g = \log\Gamma$
dan $x = \frac12$ memberikan $\sqrt{n-1} \leq
\frac{\Gamma(n+\frac12)}{\Gamma(n)} \leq \sqrt n$, sehingga
nisbahnya terhadap $\sqrt n$ terjepit antara $\sqrt{1 - \frac1n}$
dan $1$.

\textbf{21.} Dari pertanyaan 20, $(2n)! =
\frac{4^nn!}{\sqrt\pi}\Gamma(n + \tfrac12)$, sehingga
\[
\binom{2n}n = \frac{(2n)!}{(n!)^2}
= \frac{4^n\,\Gamma(n+\frac12)}{\sqrt\pi\;n!}
= \frac{4^n}{\sqrt\pi}\cdot
\frac{\Gamma(n+\frac12)}{n\,\Gamma(n)}
\sim \frac{4^n}{\sqrt\pi}\cdot\frac{\sqrt n}{n}
= \frac{4^n}{\sqrt{\pi n}} .
\]
Secara numerik, $4^{10}/\sqrt{10\pi} = 1048576/5.6050 \approx
187078.6$, terhadap $\binom{20}{10} = 184756$: yakni nisbah $0.9876$
--- sehingga galatnya $O(1/n)$, yang terlihat di $n = 10$.

\textbf{22.} Kesetaraannya: $\Gamma(\frac12) = \sqrt\pi
\Leftrightarrow G = \sqrt\pi$ (lewat substitusi $x = u^2$ pada
pertanyaan 3) $\Leftrightarrow B(\frac12, \frac12) = \pi$
(lewat rumus Euler); sedangkan penggandaan di $z = n$ \emph{adalah}
bentuk tertutup $\Gamma(n + \frac12)$, yang \emph{adalah}
taksiran binomial pusatnya, sampai lema kemiringannya. Teknik
parameternya (Bagian I--II) menghitung limit besaran yang
bergerak dan tak tergantikan bila sebuah perubahan bentuk yang sejati
hadir (sebab integral Dirichlet tak mempunyai bukti kecembungan); sedangkan
teknik kecembungannya tak menghitung apa pun, melainkan melarang semua hal
--- ia unggul pada ketunggalan dan persamaan fungsional
(Gauss, Euler, Legendre dalam tiga sapuan), tempat
pendiferensialan akan tenggelam dalam perhitungan. Jadi seorang
analis yang \omterm{def:b3:complete:complete}{lengkap} membawa keduanya.

\textbf{23.} Dengan $t = \sin^2\theta$, $\dd t =
2\sin\theta\cos\theta\,\dd\theta =
2\,t^{1/2}(1-t)^{1/2}\,\dd\theta$, sehingga
\[
W_p = \int_0^1 t^{p/2}\,
\frac{\dd t}{2\,t^{1/2}(1-t)^{1/2}}
= \frac12\int_0^1 t^{\frac{p+1}2 - 1}(1-t)^{\frac12 - 1}\dd t
= \frac12\,B\Bigl(\frac{p+1}2, \frac12\Bigr).
\]
Lalu rekursi Betanya dengan $x = \frac{n+1}2$ dan $y = \frac12$
memberikan
\[
W_{n+2} = \frac12\,
\frac{(n+1)/2}{(n+2)/2}\,B\Bigl(\frac{n+1}2, \frac12\Bigr)
= \frac{n+1}{n+2}\,W_n .
\]
Mulai dari $W_0 = \frac\pi2$ dan $W_1 = 1$:
\[
W_{2n} = \frac\pi2\prod_{k=1}^n\frac{2k-1}{2k},
\qquad
W_{2n+1} = \prod_{k=1}^n\frac{2k}{2k+1} .
\]
Karena $\sin^{n+1} \leq \sin^n$ pada $\intcc0{\pi/2}$, maka
barisan $(W_n)$ tak naik, sehingga
\[
1 \geq \frac{W_{2n+1}}{W_{2n}} \geq
\frac{W_{2n+1}}{W_{2n-1}} = \frac{2n}{2n+1}
\longrightarrow 1 .
\]
Padahal bentuk tertutupnya memberikan
\[
\frac{W_{2n+1}}{W_{2n}} = \frac2\pi
\prod_{k=1}^n\frac{(2k)^2}{(2k-1)(2k+1)}
= \frac2\pi\prod_{k=1}^n\frac{4k^2}{4k^2-1},
\]
dan membiarkan $n \to \infty$ melahirkan hasil kali Wallis. (Lewat
rumus Euler, $W_p = \frac{\sqrt\pi}2\,
\Gamma(\frac{p+1}2)/\Gamma(\frac p2 + 1)$: jadi Wallis adalah
integral Gauss dalam kostum yang lain lagi.)

\textbf{24.} Turunan-$b$ integrannya adalah
$-2x\,\eu^{-x^2}\sin(2bx)$, yang terdominasi oleh $2\abs
x\,\eu^{-x^2} \in L^1(\R)$ secara seragam terhadap $b$: sehingga $\Phi$ bersifat
$\mathcal C^1$ dengan
\[
\Phi'(b) = -\int_{-\infty}^{+\infty}
2x\,\eu^{-x^2}\sin(2bx)\,\dd x .
\]
Lalu dengan pengintegralan parsial memakai $u = \sin(2bx)$ dan $\dd v =
-2x\,\eu^{-x^2}\dd x$ (sehingga $v = \eu^{-x^2}$), suku batasnya
lenyap dan
\[
\Phi'(b) = -2b\int_{-\infty}^{+\infty}
\eu^{-x^2}\cos(2bx)\,\dd x = -2b\,\Phi(b).
\]
Karena itu $\bigl(\Phi(b)\,\eu^{b^2}\bigr)' = 0$ dan $\Phi(b) =
\Phi(0)\,\eu^{-b^2} = \sqrt\pi\,\eu^{-b^2}$ menurut Bagian I. Sampai
penormalannya, ini mengatakan bahwa transformasi Fourier
$\eu^{-x^2}$ kembali berupa fungsi Gauss --- yakni titik tetap yang menjadi
poros teori pembalikan \cref{ch:b3:fouriertransform}.

\textbf{25.} Untuk $0 < a < b$ dan $x > 0$,
\[
0 \leq \frac{\eu^{-ax} - \eu^{-bx}}{x}
= \int_a^b \eu^{-sx}\,\dd s \leq (b - a)\,\eu^{-ax},
\]
sehingga integral $I(a)$ konvergen (secara Lebesgue); dan batas titik demi
titiknya juga menunjukkan bahwa integrannya diperluas secara \omterm{def:b3:topology:continuity}{kontinu} oleh $b -
a$ di $x = 0$. Tetapkan $b$; lalu pada $\intco{a_0}{+\infty}$
turunan-$a$ integrannya adalah $-\eu^{-ax}$, yang terdominasi oleh
$\eu^{-a_0x} \in L^1(\intoo0{+\infty})$, sehingga $I$ bersifat $\mathcal
C^1$ pada $\intoo0b$ dengan
\[
I'(a) = -\int_0^{+\infty}\eu^{-ax}\dd x = -\frac1a,
\qquad\text{sehingga}\qquad
I(a) = -\log a + c .
\]
Batas dua sisinya memberikan $0 \leq I(a) \leq (b-a)/a \to 0$ saat
$a \to b^-$, sehingga $c = \log b$ dan $I(a) = \log\frac ba$. Sedangkan
kesingularan terhapuskannya diperlukan \emph{di $0$}: sebab masing-masing suku
$\eu^{-ax}/x$ secara tersendiri berintegral divergen (secara
logaritmik) di dekat $0$, dan hanya peniadaan orde pertamanya
$\eu^{-ax} - \eu^{-bx} = O(x)$ yang membuat selisihnya
\omterm{def:b3:lebesgue:l1}{terintegralkan} di sana; sedangkan di tak hingga masing-masing suku tak berbahaya
dengan sendirinya.

\end{solution}
